\mxnewpage\section*{EV1-T1 -- DS commun}\label{EV1-T1}
\addcontentsline{toc}{section}{EV1-T1 -- DS commun}
\textbf{Sujet d’entraînement}. Chapitres 1 à 4.\par Durée : 2 h. Sans calculatrice. Note sur 20. 14 novembre 2026.
\subsection*{Partie 1 -- Automatismes (6 points)}
Douze questions, \num{0,5} point chacune. Une seule réponse correcte ; aucune justification demandée. Une réponse incorrecte ou absente rapporte 0, sans retrait de points.
\begin{enumerate}
\item Calculer $3+2\times3$.\par A. $12$\quad B. $9$\quad C. $8$\quad D. $15$
\item Calculer $\frac12+\frac13$.\par A. $\frac{5}{3}$\quad B. $\frac{2}{5}$\quad C. $\frac{5}{6}$\quad D. $\frac{1}{6}$
\item Calculer $\dfrac{2^3}{2^{1}}$.\par A. $4$\quad B. $8$\quad C. $2$\quad D. $1$
\item Une valeur de 50 augmente de \pc{20}. Valeur finale ?\par A. $10$\quad B. $40$\quad C. $60$\quad D. $70$
\item Une valeur de 80 est obtenue après une baisse de \pc{20}. Valeur initiale ?\par A. $60$\quad B. $64$\quad C. $96$\quad D. $100$
\item Hausse de \pc{25} puis baisse de \pc{20} : coefficient global ?\par A. $\frac{4}{5}$\quad B. $\frac{21}{20}$\quad C. $\frac{9}{10}$\quad D. $1$
\item La solution de $2x+3=0$ vaut\par A. $\frac{3}{2}$\quad B. $\frac{-3}{2}$\quad C. $6$\quad D. $-6$
\item La pente entre $\pt{0;3}$ et $\pt{2;9}$ vaut\par A. $2$\quad B. $-3$\quad C. $6$\quad D. $3$
\item La moyenne de $3, 5, 7$ vaut\par A. $6$\quad B. $15$\quad C. $3$\quad D. $5$
\item Un secteur de $90^\circ$ représente quelle fraction du total ?\par A. $\frac{1}{2}$\quad B. $\frac{1}{4}$\quad C. $\frac{3}{4}$\quad D. $\frac{1}{3}$
\item Une probabilité vaut $\dfrac{2}{5}$. Son complément vaut\par A. $\frac{2}{5}$\quad B. $\frac{3}{5}$\quad C. $\frac{5}{2}$\quad D. $\frac{1}{5}$
\item Pour $f(x)=x^2+3$, $f(-2)$ vaut\par A. $7$\quad B. $5$\quad C. $12$\quad D. $-1$
\end{enumerate}
\mxnewpage\subsection*{Partie 2 -- Exercices indépendants (14 points)}
Toute démarche doit être justifiée. Les résultats exacts sont attendus sauf indication contraire.
\mxneed{12}
\subsection*{Exercice 1 -- Un fonds de livres (5 points)}
Une bibliothèque possède 120 livres à la fin du mois de référence, noté mois 0. À la fin de chaque mois suivant, elle ajoute 18 livres à son fonds ; aucun livre n’en est retiré. On note $u_n$ le nombre de livres disponibles à la fin du mois $n$, pour tout entier naturel $n$.
\begin{enumerate}
\item Déterminer la nature de $(u_n)$, sa raison et son expression en fonction de $n$. \hfill(1 pt)
\item Calculer $u_8$, puis $u_0+u_1+\cdots+u_8$. Expliquer ce que cette somme cumule et pourquoi elle ne compte pas des livres tous distincts. \hfill(2 pts)
\item Déterminer le premier mois où le fonds dépasse strictement 300 livres. Contrôler le mois précédent. \hfill(2 pts)
\end{enumerate}
\mxneed{12}
\subsection*{Exercice 2 -- Aire et choix de forme (5 points)}
Un rectangle a un périmètre de 28 m. Si l’un de ses côtés mesure $x$ mètres, l’autre mesure $14-x$ mètres. Son aire est modélisée par la fonction $A$ définie sur $\intff{0}{14}$ par $A(x)=x(14-x)$, en mètres carrés. Les bornes représentent les rectangles aplatis.
\begin{enumerate}
\item Développer puis écrire A sous forme canonique. \hfill(1 pt)
\item Déterminer l’aire maximale et la valeur de x correspondante sans dérivée. \hfill(2 pts)
\item Résoudre $A(x)\geq40$ sur le domaine. \hfill(2 pts)
\end{enumerate}
\mxneed{12}
\subsection*{Exercice 3 -- Un angle droit (4 points)}
Le plan est muni d’un repère orthonormé. On donne les points $A\pt{0;1}$, $B\pt{3;2}$ et $C\pt{-1;4}$. Les longueurs sont exprimées dans l’unité du repère et les aires en unités d’aire.
\begin{enumerate}
\item Calculer les coordonnées des vecteurs $\vect{AB}$ et $\vect{AC}$. \hfill(1 pt)
\item Démontrer que ces vecteurs sont orthogonaux et préciser la nature du triangle. \hfill(1 pt)
\item Calculer la longueur $BC$ puis l’aire du triangle $ABC$. \hfill(2 pts)
\end{enumerate}
